% در این فایل، دستورها و تنظیمات مورد نیاز، آورده شده است.
%-------------------------------------------------------------------------------------------------------------------

% در ورژن جدید زی‌پرشین برای تایپ متن‌های ریاضی، این سه بسته، حتماً باید فراخوانی شود
\usepackage{amsthm,amssymb,amsmath}
% بسته‌ای برای تنطیم حاشیه‌های بالا، پایین، چپ و راست صفحه
\usepackage[top=30mm, bottom=25mm, left=25mm, right=35mm]{geometry}
% بسته‌‌ای برای ظاهر شدن شکل‌ها و تصاویر متن
\usepackage{graphicx}
% بسته‌ای برای رسم کادر
\usepackage{framed} 
% بسته‌‌ای برای چاپ شدن خودکار تعداد صفحات در صفحه «معرفی پایان‌نامه»
\usepackage{lastpage}
% بسته‌ و دستوراتی برای ایجاد لینک‌های رنگی با امکان جهش
% \usepackage[pagebackref=false,colorlinks,linkcolor=black,citecolor=black]{hyperref}
% چنانچه قصد پرینت گرفتن نوشته خود را دارید، خط بالا را غیرفعال و  از دستور زیر استفاده کنید چون در صورت استفاده از دستور زیر‌‌، 
% لینک‌ها به رنگ سیاه ظاهر خواهند شد که برای پرینت گرفتن، مناسب‌تر است
% \usepackage[pagebackref=false]{hyperref}
% بسته‌ لازم برای تنظیم سربرگ‌ها
\usepackage{fancyhdr}
\setlength{\headheight}{16pt}
%
\usepackage{setspace}
\usepackage{algorithm}
\usepackage{algorithmic}
\usepackage{subfigure}
\usepackage{titletoc}


% بسته‌ای برای ظاهر شدن «مراجع» و «نمایه» در فهرست مطالب
\usepackage[nottoc]{tocbibind}
% دستورات مربوط به ایجاد نمایه
\usepackage{makeidx}
\makeindex
% تعریف شیم‌های مورد نیاز برای نسخه‌های جدید بسته bidi در برخی محیط‌ها
% Removal of redundant and potentially broken definitions
\makeatletter
% شیم‌های سازگاری برای نسخه‌های جدید لاتک و بسته‌های bidi/xepersian
\providecommand{\IfClassLoadedT}[2]{\@ifclassloaded{#1}{#2}{}}
\providecommand{\IfClassLoadedF}[2]{\@ifclassloaded{#1}{}{#2}}
\providecommand{\IfClassLoadedTF}[3]{\@ifclassloaded{#1}{#2}{#3}}
\providecommand{\IfPackageLoadedT}[2]{\@ifpackageloaded{#1}{#2}{}}
\providecommand{\IfPackageLoadedF}[2]{\@ifpackageloaded{#1}{}{#2}}
\providecommand{\IfPackageLoadedTF}[3]{\@ifpackageloaded{#1}{#2}{#3}}
\providecommand{\@iffileloaded}[3]{\expandafter\ifx\csname ver@#1\endcsname\relax#3\else#2\fi}
\providecommand{\IfFileLoadedT}[2]{\@iffileloaded{#1}{#2}{}}
\providecommand{\IfFileLoadedF}[2]{\@iffileloaded{#1}{}{#2}}
\providecommand{\IfFileLoadedTF}[3]{\@iffileloaded{#1}{#2}{#3}}
\makeatother

%%%%%%%%%%%%%%%%%%%%%%%%%%
% فراخوانی بسته زی‌پرشین و تعریف قلم فارسی و انگلیسی
\usepackage[pagebackref=false,colorlinks,linkcolor=black,citecolor=black]{hyperref}
\usepackage{notoccite}

\ExplSyntaxOn
\cs_if_exist:NTF \g__hyp_linknestlevel_int { } { \int_new:N \g__hyp_linknestlevel_int }
\cs_if_exist:NTF \l__hyp_annot_GoTo_bool { } { \bool_new:N \l__hyp_annot_GoTo_bool }
\cs_if_exist:NTF \l__hyp_annot_URI_bool { } { \bool_new:N \l__hyp_annot_URI_bool }
\cs_if_exist:NTF \l__hyp_href_url_ismap_bool { } { \bool_new:N \l__hyp_href_url_ismap_bool }
\cs_if_exist:NTF \l__hyp_text_enc_uri_print_tl { } { \tl_new:N \l__hyp_text_enc_uri_print_tl }
\cs_if_exist:NTF \l__hyp_uri_tmpa_tl { } { \tl_new:N \l__hyp_uri_tmpa_tl }

% Defining missing l3pdf/pdfdict commands if not present
\cs_if_exist:NTF \pdfdict_put:nno { } { \cs_new:Npn \pdfdict_put:nno #1#2#3 { } }
\cs_if_exist:NTF \pdfdict_put:nnn { } { \cs_new:Npn \pdfdict_put:nnn #1#2#3 { } }
\cs_if_exist:NTF \pdfdict_use:n { } { \cs_new:Npn \pdfdict_use:n #1 { } }
\cs_if_exist:NTF \pdfannot_dict_put:nne { } { \cs_new:Npn \pdfannot_dict_put:nne #1#2#3 { } }
\cs_if_exist:NTF \pdfannot_link:nen { } { \cs_new:Npn \pdfannot_link:nen #1#2#3 { } }
\cs_if_exist:NTF \__hyp_check_link_nesting:TF { } { \cs_new:Npn \__hyp_check_link_nesting:TF #1#2 { #1 } }
\cs_if_exist:NTF \__hyp_text_pdfstring:eoN { } { \cs_new:Npn \__hyp_text_pdfstring:eoN #1#2#3 { } }
\ExplSyntaxOff

% تعریف دستورات گمشده برای جلوگیری از خطای renewcommand در بسته‌های داخلی
\makeatletter
\providecommand{\foottextfont}{}
\providecommand{\LTRfoottextfont}{}
\providecommand{\RTLfoottextfont}{}
\providecommand{\abstractname}{}
\providecommand{\latinabstract}{}
\providecommand{\AutoPersianDigits}{}
\makeatother

\usepackage{xepersian}
\settextfont[Path=font/,Extension=.ttf,Scale=1,
    BoldFont=XB NiloofarBd,
    ItalicFont=XB NiloofarIt,
    BoldItalicFont=XB NiloofarBdIt
]{XB Niloofar}
\setlatintextfont[Scale=0.9]{Times New Roman}

%%%%%%%%%%%%%%%%%%%%%%%%%%
% چنانچه می‌خواهید اعداد در فرمول‌ها، انگلیسی باشد، خط زیر را غیرفعال کنید
\setdigitfont[Path=font/,Extension=.ttf,Scale=1,
    BoldFont=XB NiloofarBd,
    ItalicFont=XB NiloofarIt,
    BoldItalicFont=XB NiloofarBdIt
]{XB Niloofar}
%%%%%%%%%%%%%%%%%%%%%%%%%%
% تعریف قلم‌های فارسی و انگلیسی اضافی برای استفاده در بعضی از قسمت‌های متن
\defpersianfont\titlefont[Path=font/,Extension=.ttf,Scale=1]{XB Titre}
% \defpersianfont\iranic[Scale=1.1]{XB Zar Oblique}%Italic}%
% \defpersianfont\nastaliq[Scale=1.2]{IranNastaliq}

%%%%%%%%%%%%%%%%%%%%%%%%%%
% دستوری برای حذف کلمه «چکیده»
\renewcommand{\abstractname}{}
% دستوری برای حذف کلمه «abstract»
%\renewcommand{\latinabstract}{}
% دستوری برای تغییر نام کلمه «اثبات» به «برهان»
\renewcommand\proofname{\textbf{برهان}}
% دستوری برای تغییر نام کلمه «کتاب‌نامه» به «مراجع»
\renewcommand{\bibname}{مراجع}
% دستوری برای تعریف واژه‌نامه انگلیسی به فارسی
\newcommand\persiangloss[2]{#1\dotfill\lr{#2}\\}
% دستوری برای تعریف واژه‌نامه فارسی به انگلیسی 
\newcommand\englishgloss[2]{#2\dotfill\lr{#1}\\}
% تعریف دستور جدید «\پ» برای خلاصه‌نویسی جهت نوشتن عبارت «پروژه/پایان‌نامه/رساله»
\newcommand{\پ}{پروژه/پایان‌نامه/رساله }

%\newcommand\BackSlash{\char`\\}

%%%%%%%%%%%%%%%%%%%%%%%%%%
\SepMark{-}

% تعریف و نحوه ظاهر شدن عنوان قضیه‌ها، تعریف‌ها، مثال‌ها و ...
\theoremstyle{definition}
\newtheorem{definition}{تعریف}[section]
\theoremstyle{plain}
\newtheorem{theorem}[definition]{قضیه}
\newtheorem{lemma}[definition]{لم}
\newtheorem{proposition}[definition]{گزاره}
\newtheorem{corollary}[definition]{نتیجه}
\newtheorem{remark}[definition]{ملاحظه}
\theoremstyle{definition}
\newtheorem{example}[definition]{مثال}

%\renewcommand{\theequation}{\thechapter-\arabic{equation}}
%\def\bibname{مراجع}
\numberwithin{algorithm}{chapter}
\def\listalgorithmname{فهرست الگوریتم‌ها}
\def\listfigurename{فهرست تصاویر}
\def\listtablename{فهرست جداول}

%%%%%%%%%%%%%%%%%%%%%%%%%%%%
% دستورهایی برای سفارشی کردن سربرگ صفحات
% \newcommand{\SetHeader}{
% \csname@twosidetrue\endcsname
% \pagestyle{fancy}
% \fancyhf{} 
% \fancyhead[OL,EL]{\thepage}
% \fancyhead[OR]{\small\rightmark}
% \fancyhead[ER]{\small\leftmark}
% \renewcommand{\chaptermark}[1]{%
% \markboth{\thechapter-\ #1}{}}
% }
%%%%%%%%%%%%5
%\def\MATtextbaseline{1.5}
%\renewcommand{\baselinestretch}{\MATtextbaseline}
\doublespacing
%%%%%%%%%%%%%%%%%%%%%%%%%%%%%
% دستوراتی برای اضافه کردن کلمه «فصل» در فهرست مطالب با استفاده از titletoc
\titlecontents{chapter}
  [5em]
  {\addvspace{1em}\bfseries}
  {\contentslabel[\chaptername~\thecontentslabel:]{5em}}
  {\hspace*{-5em}}
  {\titlerule*[0.5pc]{.}\contentspage}

\titlecontents{section}
  [5em]
  {}
  {\contentslabel{2em}}
  {\hspace*{-2em}}
  {\titlerule*[0.5pc]{.}\contentspage}

\titlecontents{subsection}
  [8em]
  {}
  {\contentslabel{3em}}
  {\hspace*{-3em}}
  {\titlerule*[0.5pc]{.}\contentspage}

\titlecontents{subsubsection}
  [12em]
  {}
  {\contentslabel{4em}}
  {\hspace*{-4em}}
  {\titlerule*[0.5pc]{.}\contentspage}

\makeatletter
{\def\thebibliography#1{\chapter*{\refname\@mkboth
   {\uppercase{\refname}}{\uppercase{\refname}}}\list
   {[\arabic{enumi}]}{\settowidth\labelwidth{[#1]}
   \rightmargin\labelwidth
   \advance\rightmargin\labelsep
   \advance\rightmargin\bibindent
   \itemindent -\bibindent
   \listparindent \itemindent
   \parsep \z@
   \usecounter{enumi}}
   \def\newblock{}
   \sloppy
   \sfcode`\.=1000\relax}}
\makeatother
